\chapter{الإعدادات}
\label{ch:settings}

\section{ماذا ستنجز؟}

يجمع تطبيق «الإعدادات» القرارات التي تسري على المؤسسة كلها: هويتها، ومن يدخل إليها، وما يصل إلى المتعاملين، والقنوات التي تتواصل عبرها. يتولى هذه الصفحات عادةً مسؤول النظام، أما الموظف فلا يحتاج إليها في عمله اليومي. ولا تظهر بعض الصفحات إلا لمن يملك صلاحية إدارتها.

\begin{rcnote}[مثال: شركة الأفق]
يراجع مسؤول النظام في شركة الأفق إعدادات المؤسسة مرة كل ثلاثة أشهر. وعندما تطلب سارة الخطيب تفعيل التحقق بخطوتين، يفتح الصفحة من مجموعة «الوصول والأمان»، ويغيّر الإعداد، ثم يتأكد من أثره قبل إبلاغ الفريق.
\end{rcnote}

\begin{figure}[H]
\centering
\begin{tikzpicture}[node distance=0.5cm]
\node[stepBox] (a) {\RL{اختر المجموعة والصفحة}};
\node[stepBox,below=of a] (b) {\RL{غيّر قيمة واحدة واحفظها}};
\node[stepBox,below=of b] (c) {\RL{افتح الصفحة التي تتأثر بها}};
\draw[flowsTo] (a) -- (b);
\draw[flowsTo] (b) -- (c);
\end{tikzpicture}
\rclegend
\caption{تغيير إعداد بأمان}
\label{fig:settings-change-flow}
\shotnote{يؤثر الإعداد في المؤسسة كلها، لا في حسابك وحده. لذلك غيّر قيمة واحدة في كل مرة، ثم افتح الصفحة التي تتأثر بها وتأكد من النتيجة قبل الانتقال إلى التغيير التالي؛ فإذا ظهرت مشكلة عرفت سببها وسهل عليك الرجوع عنها.}
\end{figure}

\section{دليل صفحات الإعدادات}
\label{sec:settings-pages}
تتبع الصفحات المجموعات الثماني التي تعرضها صفحة «الإعدادات»، وبالترتيب نفسه. ابدأ بمجموعتي «المؤسسة» و«الوصول والأمان»، ثم انتقل إلى المجموعات التي تستخدمها مؤسستك؛ فإذا لم تكن لديكم بوابة عامة مثلاً، فتجاوز صفحاتها.

\begin{table}[H]
\centering
\caption{المؤسسة}
\label{tab:settings-pages-organization}
\begin{tabularx}{\linewidth}{R{4.2cm} Y}
\rcHeaderRow
\thh{الصفحة وقسمها} & \thh{ماذا تقدم هذه الصفحة؟}\\
إعدادات المؤسسة — المؤسسة & تحديث هوية المؤسسة وإعداداتها الافتراضية.\\
إعدادات الإشعارات الافتراضية — المؤسسة & ضبط تفضيلات الإشعارات على مستوى المؤسسة.\\
مهام التصدير — المؤسسة & مراجعة ملفات تصدير البيانات وتنزيلها.\\
التقارير — المؤسسة & الفترات وعرض العملة والإرسال والمفاتيح التي يبدأ بها كل تقرير.\\
\end{tabularx}
\end{table}

\begin{table}[H]
\centering
\caption{الوصول والأمان}
\label{tab:settings-pages-access}
\begin{tabularx}{\linewidth}{R{4.2cm} Y}
\rcHeaderRow
\thh{الصفحة وقسمها} & \thh{ماذا تقدم هذه الصفحة؟}\\
الأدوار — الوصول والأمان & تحديد ما يراه كل عضو وما يفعله.\\
التحقق بخطوتين — الوصول والأمان & اشتراط خطوة ثانية لتسجيل الدخول ومتابعة من فعّلها.\\
الجلسات والأجهزة — الوصول والأمان & مدة الجلسة وتسجيل الخروج عند عدم النشاط وتنبيهات الأجهزة الجديدة.\\
مفاتيح الأعضاء — الوصول والأمان & إدارة مفاتيح التوقيع المسجلة للأعضاء.\\
إصدار رمز تسجيل — الوصول والأمان & إصدار رموز تسجيل مفاتيح التوقيع للأعضاء.\\
سجل التدقيق — الوصول والأمان & مراجعة أحداث تسجيل مفاتيح الأعضاء واستخدامها.\\
الواجهة البرمجية والتكاملات — الوصول والأمان & إدارة رموز \en{API} وقواعد اتصال الأنظمة الأخرى بالمؤسسة.\\
\end{tabularx}
\end{table}

\begin{table}[H]
\centering
\caption{البوابة العامة}
\label{tab:settings-pages-portal}
\begin{tabularx}{\linewidth}{R{4.2cm} Y}
\rcHeaderRow
\thh{الصفحة وقسمها} & \thh{ماذا تقدم هذه الصفحة؟}\\
البوابة العامة — البوابة العامة & إعداد البوابة العامة للمؤسسة.\\
محتوى البوابة — البوابة العامة & الفئات المستهدفة والصفحات والأسئلة الشائعة والإعلانات.\\
حسابات المتعاملين — البوابة العامة & اختيار بيانات الملف الشخصي وطريقة التحقق منها.\\
الرسائل — البوابة العامة & ما يُبلَّغ به المتعاملون عن طلباتهم وعبر أي قناة.\\
\end{tabularx}
\end{table}

\begin{table}[H]
\centering
\caption{النماذج والنشر}
\label{tab:settings-pages-forms}
\begin{tabularx}{\linewidth}{R{4.2cm} Y}
\rcHeaderRow
\thh{الصفحة وقسمها} & \thh{ماذا تقدم هذه الصفحة؟}\\
النشر والاعتماد — النماذج والنشر & التغييرات التي تحتاج إلى اعتماد قبل تطبيقها، ومن يعتمدها.\\
التكليف بالنماذج — النماذج والنشر & التذكير والتصعيد وأيام العمل والحد المعتاد لكل شخص.\\
تصنيفات الفحوص — النماذج والنشر & تصنيفات تُجمَّع بها فحوص الجودة للمراجعين ولوحات المعلومات.\\
\end{tabularx}
\end{table}

\begin{table}[H]
\centering
\caption{قنوات التواصل}
\label{tab:settings-pages-channels}
\begin{tabularx}{\linewidth}{R{4.2cm} Y}
\rcHeaderRow
\thh{الصفحة وقسمها} & \thh{ماذا تقدم هذه الصفحة؟}\\
نطاقات البريد الإلكتروني — قنوات التواصل & إدارة النطاقات الموثقة لبريد المؤسسة.\\
توقيعات البريد الإلكتروني — قنوات التواصل & تصميم توقيعات بريد المؤسسة وتعيينها.\\
نطاقات الروابط القصيرة — قنوات التواصل & إدارة النطاقات المخصصة للروابط المختصرة.\\
حسابات واتساب — قنوات التواصل & ربط حسابات واتساب وإدارتها.\\
\end{tabularx}
\end{table}

\begin{table}[H]
\centering
\caption{الذكاء الاصطناعي}
\label{tab:settings-pages-ai}
\begin{tabularx}{\linewidth}{R{4.2cm} Y}
\rcHeaderRow
\thh{الصفحة وقسمها} & \thh{ماذا تقدم هذه الصفحة؟}\\
توفر الذكاء الاصطناعي — الذكاء الاصطناعي & اختيار أين يتوفر الذكاء الاصطناعي، للجميع أو لأعضاء محددين.\\
حدود الذكاء الاصطناعي — الذكاء الاصطناعي & تحديد حدود استخدام مزايا الذكاء الاصطناعي في المؤسسة.\\
\end{tabularx}
\end{table}

\begin{table}[H]
\centering
\caption{اللغات والعملات والوحدات}
\label{tab:settings-pages-languages}
\begin{tabularx}{\linewidth}{R{4.2cm} Y}
\rcHeaderRow
\thh{الصفحة وقسمها} & \thh{ماذا تقدم هذه الصفحة؟}\\
اللغات — اللغات والعملات والوحدات & إدارة لغات المؤسسة.\\
العملات — اللغات والعملات والوحدات & إدارة عملات المؤسسة.\\
وحدات القياس — اللغات والعملات والوحدات & الوحدات التي تظهر في الأسئلة والحقول، وتحويلاتها.\\
\end{tabularx}
\end{table}

\begin{table}[H]
\centering
\caption{الإدارة}
\label{tab:settings-pages-management}
\begin{tabularx}{\linewidth}{R{4.2cm} Y}
\rcHeaderRow
\thh{الصفحة وقسمها} & \thh{ماذا تقدم هذه الصفحة؟}\\
مسارات الموافقة — الإدارة & مسارات موافقة محفوظة بأسماء وإصدارات، تُستخدم في أكثر من مكان.\\
قوالب التقييم — الإدارة & بطاقات تقييم قابلة لإعادة الاستخدام: معايير ودرجات وفئات.\\
المتابعة — الإدارة & أسباب التعطل وغيرها من الكلمات المشتركة في لوحات المتابعة.\\
حفظ السجلات — الإدارة & قواعد الاحتفاظ والتجميد ومراجعات الإتلاف والسلامة.\\
\end{tabularx}
\end{table}
